% ═══════════════════════════════════════════════════════════════
% LEHRERHINWEISE: Repaso Unidad 5 + Números 100-1000
% Spanisch Klasse 8 | 90 Min. (Doppelstunde)
% ═══════════════════════════════════════════════════════════════

\documentclass[11pt, a4paper]{article}

\newif\ifswmodus
\swmodusfalse

\usepackage[utf8]{inputenc}
\usepackage[T1]{fontenc}
\usepackage[ngerman]{babel}
\usepackage{helvet}
\renewcommand{\familydefault}{\sfdefault}

\usepackage[margin=2cm]{geometry}
\usepackage{enumitem}
\usepackage{tabularx}
\usepackage{booktabs}
\usepackage{longtable}

\usepackage{xcolor}
\usepackage{amssymb}
\usepackage{tcolorbox}
\tcbuselibrary{skins}

\usepackage{fancyhdr}
\usepackage{lastpage}
\usepackage{needspace}
\usepackage[hyphens,spaces,obeyspaces]{url}
\usepackage[colorlinks=true,linkcolor=spanisch-color,urlcolor=spanisch-color,breaklinks=true]{hyperref}

% ═══════════════════════════════════════════════════════════════
% FARBEN
% ═══════════════════════════════════════════════════════════════

\definecolor{spanisch-color}{HTML}{c41e3a}
\definecolor{grau-color}{HTML}{888888}
\definecolor{hellgrau-color}{HTML}{f5f5f5}
\definecolor{basis-color}{HTML}{2a7a4b}
\definecolor{standard-color}{HTML}{2c5aa0}
\definecolor{erweiterung-color}{HTML}{7b1fa2}
\definecolor{loesung-color}{HTML}{c62828}
\definecolor{orange-color}{HTML}{b35c00}

\definecolor{sw-dark}{HTML}{000000}
\definecolor{sw-gray}{HTML}{666666}
\definecolor{sw-white}{HTML}{ffffff}

\ifswmodus
    \colorlet{spanisch}{sw-dark}
    \colorlet{grau}{sw-gray}
    \colorlet{hellgrau}{sw-white}
    \colorlet{basis}{sw-dark}
    \colorlet{standard}{sw-dark}
    \colorlet{erweiterung}{sw-dark}
    \colorlet{loesung}{sw-dark}
    \colorlet{orange}{sw-dark}
\else
    \colorlet{spanisch}{spanisch-color}
    \colorlet{grau}{grau-color}
    \colorlet{hellgrau}{hellgrau-color}
    \colorlet{basis}{basis-color}
    \colorlet{standard}{standard-color}
    \colorlet{erweiterung}{erweiterung-color}
    \colorlet{loesung}{loesung-color}
    \colorlet{orange}{orange-color}
\fi

\newcommand{\fachfarbe}{spanisch}

% ═══════════════════════════════════════════════════════════════
% VARIABLEN
% ═══════════════════════════════════════════════════════════════

\newcommand{\thema}{Repaso Unidad 5 + Números 100-1000}
\newcommand{\fach}{Spanisch}
\newcommand{\klasse}{8}
\newcommand{\dauer}{90 Min. (Doppelstunde)}
\newcommand{\lerngruppengroesse}{24}
\newcommand{\datum}{\today}

% ═══════════════════════════════════════════════════════════════
% KOPF- UND FUSSZEILE
% ═══════════════════════════════════════════════════════════════

\pagestyle{fancy}
\fancyhf{}
\renewcommand{\headrulewidth}{0.4pt}
\renewcommand{\footrulewidth}{0.4pt}

\lhead{\fach{} -- Klasse \klasse}
\chead{\textbf{Lehrerhinweise}}
\rhead{\datum}

\lfoot{}
\cfoot{}
\rfoot{Seite \thepage\ von \pageref{LastPage}}

% ═══════════════════════════════════════════════════════════════
% BEFEHLE
% ═══════════════════════════════════════════════════════════════

\newcommand{\B}{\textcolor{basis}{\textbf{[B]}}}
\newcommand{\Std}{\textcolor{standard}{\textbf{[S]}}}
\newcommand{\E}{\textcolor{erweiterung}{\textbf{[E]}}}

\newcommand{\loesung}[1]{\textcolor{loesung}{\textbf{#1}}}
\newcommand{\es}[1]{\textcolor{\fachfarbe}{\textbf{#1}}}
\newcommand{\en}[1]{{\scriptsize\textcolor{grau}{#1}}}

\newtcolorbox{kurzplan}{
    colback=hellgrau,
    colframe=\fachfarbe,
    colbacktitle=white,
    coltitle=\fachfarbe,
    boxrule=1pt,
    arc=0pt,
    fonttitle=\bfseries,
    attach boxed title to top left={yshift=-2mm, xshift=5mm},
    boxed title style={boxrule=0pt, colback=white},
    title=Kurzplan
}

\newtcolorbox{fehlerbox}{
    colback=white,
    colframe=loesung,
    colbacktitle=white,
    coltitle=loesung,
    boxrule=0.5pt,
    arc=0pt,
    fonttitle=\bfseries,
    attach boxed title to top left={yshift=-2mm, xshift=5mm},
    boxed title style={boxrule=0pt, colback=white},
    title=Häufige Fehler
}

\newtcolorbox{materialbox}{
    colback=white,
    colframe=\fachfarbe,
    colbacktitle=white,
    coltitle=\fachfarbe,
    boxrule=0.5pt,
    arc=0pt,
    fonttitle=\bfseries,
    attach boxed title to top left={yshift=-2mm, xshift=5mm},
    boxed title style={boxrule=0pt, colback=white},
    title=Material-Checkliste
}

\newtcolorbox{digitalbox}{
    colback=white,
    colframe=orange,
    colbacktitle=white,
    coltitle=orange,
    boxrule=0.5pt,
    arc=0pt,
    fonttitle=\bfseries,
    attach boxed title to top left={yshift=-2mm, xshift=5mm},
    boxed title style={boxrule=0pt, colback=white},
    title=Digitale Materialien (optional)
}

% ═══════════════════════════════════════════════════════════════
% DOKUMENT
% ═══════════════════════════════════════════════════════════════

\begin{document}

% ═══════════════════════════════════════════════════════════════
% KOPFBEREICH
% ═══════════════════════════════════════════════════════════════

\begin{center}
    {\Large\bfseries\textcolor{\fachfarbe}{Lehrerhinweise: \thema}}\\[0.3em]
    {\normalsize \dauer{} | Klasse \klasse}
\end{center}

\vspace{10pt}

% ═══════════════════════════════════════════════════════════════
% 1. ÜBERSICHT
% ═══════════════════════════════════════════════════════════════

\section*{1. Übersicht}

\textbf{Themen:}
\begin{itemize}[nosep]
    \item \textbf{Wiederholung (Unidad 5):} Kleidung (la ropa), Farben (los colores), Modalverben (querer, poder, tener que)
    \item \textbf{Neu:} Zahlen 100--1000 (los números 100--1000)
\end{itemize}

\textbf{Lernziele:}
\begin{itemize}[nosep]
    \item SuS können Kleidungsstücke und Farben korrekt benennen und kombinieren
    \item SuS können Modalverben im Kontext konjugieren und anwenden
    \item SuS können Zahlen 100--1000 verstehen und produzieren
    \item SuS können einen Einkaufsdialog führen
\end{itemize}

% ═══════════════════════════════════════════════════════════════
% 2. MATERIAL-CHECKLISTE
% ═══════════════════════════════════════════════════════════════

\section*{2. Material-Checkliste}

\begin{materialbox}
\textbf{Pflicht:}
\begin{itemize}[label=$\square$, nosep]
    \item Lehrwerk: \textit{¿Qué pasa?} Bd. 2, S. 78--85 (Unidad 5)
    \item AB-repaso-numeros.pdf (\lerngruppengroesse{} Kopien)
    \item ML-repaso-numeros.pdf (1× für Lehrkraft)
    \item Tafel / Whiteboard
\end{itemize}
\end{materialbox}

\vspace{0.5em}

\begin{digitalbox}
\textbf{Online-Materialien (ergänzend):}

\vspace{0.5em}
\textbf{Lernpfad:}\\
\url{https://devbox42.github.io/sciencesim/spanisch/kl08/unidad5-repaso-numeros/LP-repaso-numeros.html}

\vspace{0.3em}
\textbf{Stundenleistung (28 BE):}\\
\url{https://devbox42.github.io/sciencesim/spanisch/kl08/unidad5-repaso-numeros/TEST-repaso-numeros.html}

\vspace{0.3em}
\textbf{Weitere Materialien (lokal):}
\begin{itemize}[label=$\square$, nosep]
    \item QR-LERNPFAD.pdf, QR-TEST.pdf (für Beamer)
    \item PLATZKARTEN.html (für Sitzplatzzuordnung)
\end{itemize}

\vspace{0.3em}
\textbf{Reset-PIN:} \textcolor{orange}{\textbf{2024}} (für Lehrer-Reset am Ende der Seite)
\end{digitalbox}

% ═══════════════════════════════════════════════════════════════
% 3. STUNDENVERLAUF
% ═══════════════════════════════════════════════════════════════

\section*{3. Stundenverlauf (90 Min.) / Lesson Plan (90 min)}

\textbf{1. Stunde (45 Min.) -- Wiederholung + Zahlen-Einführung}\\
\en{1st Period (45 min) -- Review + Introduction of Numbers}

\begin{tabularx}{\textwidth}{|c|l|X|l|}
\hline
\textbf{Zeit} & \textbf{Phase} & \textbf{Inhalt} & \textbf{Material} \\
\hline
0--5 & Einstieg \newline \en{Warm-up} & \textbf{Schätz-Spiel:} 3 Kleidungsstücke zeigen (Bild/real), SuS schätzen Preise auf Spanisch. \es{¿Cuánto cuesta?} \newline \en{Guessing game: Show 3 clothing items, students guess prices in Spanish.} & Bilder/Realia \\
\hline
5--15 & Reaktivierung \newline \en{Activation} & Buch S. 78--79: Kleidung + Farben wiederholen, Bilder beschreiben lassen \newline \en{Book p. 78--79: Review clothing + colors, describe images} & Buch \\
\hline
15--25 & Übung 1 \newline \en{Practice 1} & AB Aufg. 1 (Zuordnung) + Aufg. 2 (Farben) \newline \en{Worksheet ex. 1 (matching) + ex. 2 (colors)} & AB \\
\hline
25--35 & Reaktivierung \newline \en{Activation} & Buch S. 82: Modalverben wiederholen, Beispielsätze sammeln \newline \en{Book p. 82: Review modal verbs, collect example sentences} & Buch, Tafel \\
\hline
35--45 & Übung 2 \newline \en{Practice 2} & AB Aufg. 3 (Modalverben konjugieren) \newline \en{Worksheet ex. 3 (conjugate modal verbs)} & AB \\
\hline
\end{tabularx}

\vspace{1em}

\textbf{2. Stunde (45 Min.) -- Zahlen + Anwendung}\\
\en{2nd Period (45 min) -- Numbers + Application}

\begin{tabularx}{\textwidth}{|c|l|X|l|}
\hline
\textbf{Zeit} & \textbf{Phase} & \textbf{Inhalt} & \textbf{Material} \\
\hline
0--10 & Input \newline \en{Input} & \textbf{Neu:} Zahlen 100--1000 einführen (Tafel), Unregelmäßigkeiten: \es{quinientos, setecientos, novecientos} \newline \en{New: Introduce numbers 100--1000 (board), irregulars: 500, 700, 900} & Tafel \\
\hline
10--20 & Übung 3 \newline \en{Practice 3} & AB Aufg. 4 (Zahlen schreiben) + Aufg. 5 (Zahlen erkennen), chorisches Sprechen \newline \en{Worksheet ex. 4 (write numbers) + ex. 5 (recognize numbers), choral speaking} & AB \\
\hline
20--30 & Anwendung \newline \en{Application} & Buch S. 84--85: Einkaufsdialog lesen und hören, Rollenspiel vorbereiten \newline \en{Book p. 84--85: Read and listen to shopping dialogue, prepare role play} & Buch, Audio \\
\hline
30--40 & Produktion \newline \en{Production} & AB Aufg. 6 (Einkaufsdialog ergänzen), Partnerarbeit: Dialog üben \newline \en{Worksheet ex. 6 (complete shopping dialogue), pair work: practice dialogue} & AB \\
\hline
40--45 & Sicherung \newline \en{Wrap-up} & Ergebnisse besprechen, Merkkasten im AB kontrollieren, HA ansagen \newline \en{Discuss results, check info box in worksheet, assign homework} & AB, Tafel \\
\hline
\end{tabularx}

% ═══════════════════════════════════════════════════════════════
% 4. ALTERNATIVE: DIGITALER VERLAUF
% ═══════════════════════════════════════════════════════════════

\section*{4. Alternative: Digitaler Verlauf / Digital Lesson Plan}

Falls digitale Geräte verfügbar, kann die Stunde auch so gestaltet werden:\\
\en{If digital devices are available, the lesson can also be structured as follows:}

\begin{tabularx}{\textwidth}{|c|l|X|l|}
\hline
\textbf{Zeit} & \textbf{Phase} & \textbf{Inhalt} & \textbf{Material} \\
\hline
0--5 & Einstieg \newline \en{Warm-up} & QR-Code zeigen, Lernpfad öffnen \newline \en{Show QR code, open learning path} & QR-LERNPFAD.pdf \\
\hline
5--40 & Selbstlernphase \newline \en{Self-study} & SuS arbeiten Lernpfad durch (Schritte 1--12) \newline \en{Students work through learning path (steps 1--12)} & LP (digital) \\
\hline
40--45 & Pause \newline \en{Break} & Bewegungspause \newline \en{Movement break} & -- \\
\hline
45--50 & Übergang \newline \en{Transition} & QR zeigen, Platzkarten verteilen, Platznummer wählen \newline \en{Show QR, distribute seat cards, select seat number} & QR-TEST.pdf, Platzkarten \\
\hline
50--80 & Stundenleistung \newline \en{Assessment} & SuS bearbeiten Online-Stundenleistung (28 BE) \newline \en{Students complete online assessment (28 pts)} & TEST (digital) \\
\hline
80--90 & Abschluss \newline \en{Wrap-up} & Ergebnisse besprechen, Fragen klären \newline \en{Discuss results, clarify questions} & -- \\
\hline
\end{tabularx}

\vspace{0.5em}
\textbf{Hinweis:} Bei digitalem Verlauf entfällt das gedruckte AB -- der Lernpfad enthält alle Übungen interaktiv.\\
\en{Note: In the digital version, the printed worksheet is not needed -- the learning path contains all exercises interactively.}

% ═══════════════════════════════════════════════════════════════
% 5. INKLUSION (LRS / DaZ)
% ═══════════════════════════════════════════════════════════════

\section*{5. Inklusion / Inclusion}

\textbf{LRS-Unterstützung:} \en{Dyslexia support:}
\begin{itemize}[nosep]
    \item Zeitzugabe: +10 Min. für AB und Stundenleistung \en{Extra time: +10 min}
    \item Vergrößerte AB-Version (A3 oder 150\%) bereitstellen \en{Enlarged worksheet (A3 or 150\%)}
    \item Mündliche Antworten bei Aufg. 6 akzeptieren \en{Accept oral answers for ex. 6}
\end{itemize}

\vspace{0.5em}

\textbf{DaZ-Wortliste (Deutsch → Spanisch):} \en{German as second language vocabulary:}

\begin{tabular}{ll@{\hspace{2cm}}ll}
T-Shirt & la camiseta & Hose & los pantalones \\
Pullover & el jersey & Schuhe & los zapatos \\
Kleid & el vestido & Rock & la falda \\
weiß & blanco/a & schwarz & negro/a \\
rot & rojo/a & blau & azul \\
Ich möchte... & Quiero... & Wieviel kostet...? & ¿Cuánto cuesta...? \\
\end{tabular}

\vspace{1em}

% ═══════════════════════════════════════════════════════════════
% 6. HÄUFIGE FEHLER
% ═══════════════════════════════════════════════════════════════

\section*{6. Häufige Fehler / Common Errors}

\begin{fehlerbox}
\begin{tabularx}{\textwidth}{|l|l|X|}
\hline
\textbf{Fehler} & \textbf{Richtig} & \textbf{Tipp} \\
\hline
*cincocientos & \es{quinientos} & 500 ist unregelmäßig (qu- statt c-) \\
\hline
*sietecientos & \es{setecientos} & 700 kürzt die Sieben (set- statt siete-) \\
\hline
*nuevecientos & \es{novecientos} & 900 kürzt die Neun (nov- statt nueve-) \\
\hline
ciento (allein) & \es{cien} & 100 allein = cien, 100+ = ciento \\
\hline
*doscientos personas & \es{doscientas personas} & Genus beachten! (fem. = -as) \\
\hline
*yo quere & \es{yo quiero} & Stammvokalwechsel e→ie beachten \\
\hline
\end{tabularx}
\end{fehlerbox}

% ═══════════════════════════════════════════════════════════════
% 6. LÖSUNGEN
% ═══════════════════════════════════════════════════════════════

\section*{7. Lösungen AB}

\textbf{Aufgabe 1: Kleidung zuordnen} (5 BE)

\begin{tabular}{ll}
T-Shirt & $\rightarrow$ \loesung{B) la camiseta} \\
Hose & $\rightarrow$ \loesung{C) los pantalones} \\
Schuhe & $\rightarrow$ \loesung{E) los zapatos} \\
Kleid & $\rightarrow$ \loesung{A) el vestido} \\
Pullover & $\rightarrow$ \loesung{D) el jersey} \\
\end{tabular}

\vspace{0.8em}

\textbf{Aufgabe 2: Farben ergänzen} (4 BE)

\begin{tabular}{ll}
a) la camiseta & \loesung{blanca} \\
b) el jersey & \loesung{negro} \\
c) los zapatos & \loesung{marrones} \\
d) la falda & \loesung{roja} \\
\end{tabular}

\vspace{0.8em}

\textbf{Aufgabe 3: Modalverben} (5 BE)

\begin{tabular}{ll}
a) Yo \loesung{quiero} comprar un jersey. \\
b) Tú no \loesung{puedes} ir a la fiesta. \\
c) Nosotros \loesung{tenemos que} estudiar. \\
d) Ellos \loesung{quieren} una chaqueta nueva. \\
e) María \loesung{puede} hablar español. \\
\end{tabular}

\vspace{0.8em}

\textbf{Aufgabe 4: Zahlen schreiben} (6 BE)

\begin{tabular}{ll}
a) 100 = & \loesung{cien} \\
b) 350 = & \loesung{trescientos cincuenta} \\
c) 500 = & \loesung{quinientos} \\
d) 725 = & \loesung{setecientos veinticinco} \\
e) 999 = & \loesung{novecientos noventa y nueve} \\
f) 1000 = & \loesung{mil} \\
\end{tabular}

\vspace{0.8em}

\textbf{Aufgabe 5: Zahlen erkennen} (4 BE)

\begin{tabular}{ll}
a) doscientos cuarenta = & \loesung{240} \\
b) quinientos ochenta y tres = & \loesung{583} \\
c) setecientos quince = & \loesung{715} \\
d) novecientos noventa y nueve = & \loesung{999} \\
\end{tabular}

\vspace{0.8em}

\textbf{Aufgabe 6: Einkaufsdialog} (4 BE)

\begin{tabular}{ll}
1) & \loesung{Quiero comprar un jersey.} \\
2) & \loesung{Quiero uno azul.} (oder: Lo quiero azul.) \\
3) & \loesung{¿Cuánto cuesta?} \\
4) & \loesung{Bien, lo compro.} (oder: Vale, me lo llevo.) \\
\end{tabular}

% ═══════════════════════════════════════════════════════════════
% 7. BEWERTUNG
% ═══════════════════════════════════════════════════════════════

\section*{8. Bewertung (28 BE)}

\begin{tabularx}{\textwidth}{|l|c|X|}
\hline
\textbf{Aufgabe} & \textbf{BE} & \textbf{Hinweis} \\
\hline
1: Kleidung zuordnen & 5 & Je 1 BE pro richtige Zuordnung \\
\hline
2: Farben ergänzen & 4 & Je 1 BE; Genus muss stimmen \\
\hline
3: Modalverben & 5 & Je 1 BE pro richtige Form \\
\hline
4: Zahlen schreiben & 6 & Je 1 BE; Rechtschreibung beachten \\
\hline
5: Zahlen erkennen & 4 & Je 1 BE pro richtige Ziffer \\
\hline
6: Einkaufsdialog & 4 & Je 1 BE pro sinnvolle Antwort \\
\hline
\textbf{Gesamt} & \textbf{28} & \\
\hline
\end{tabularx}

\vspace{0.8em}

\textbf{Notenspiegel (14-NP, Sek I):}

\begin{tabular}{|c|c|c|c|c|c|c|c|c|c|c|c|c|c|c|}
\hline
\textbf{NP} & 14 & 13 & 12 & 11 & 10 & 9 & 8 & 7 & 6 & 5 & 4 & 3 & 2 & 1 \\
\hline
\textbf{\%} & 100 & 96 & 91 & 86 & 80 & 73 & 67 & 60 & 53 & 47 & 40 & 33 & 27 & 20 \\
\hline
\textbf{BE} & 28 & 27 & 25 & 24 & 22 & 21 & 19 & 17 & 15 & 13 & 11 & 9 & 7 & 6 \\
\hline
\end{tabular}

% ═══════════════════════════════════════════════════════════════
% 8. HAUSAUFGABE
% ═══════════════════════════════════════════════════════════════

\section*{9. Hausaufgabe}

\begin{itemize}[nosep]
    \item Lehrwerk S. 85, Aufg. 3a--c (Zahlen üben)
    \item Vokabeln lernen: Zahlen 100--1000
    \item Optional: Lernpfad LP-repaso-numeros.html (digital)
\end{itemize}

% ═══════════════════════════════════════════════════════════════
% 9. REFLEXIONSNOTIZEN
% ═══════════════════════════════════════════════════════════════

\section*{10. Reflexionsnotizen (nach Durchführung)}

\begin{tabularx}{\textwidth}{|l|X|}
\hline
\textbf{Was lief gut?} & \\[2em]
\hline
\textbf{Was lief nicht?} & \\[2em]
\hline
\textbf{Zeitmanagement} & \\[2em]
\hline
\textbf{Für nächstes Mal} & \\[2em]
\hline
\end{tabularx}

\end{document}
